%% kapa: exam=1 round=32 year=2021 period=2 form=A subject=appraisal_law subject_ko=감정평가관계법규 seq=5 number=5
%% source: Q-Net 감정평가사 기출문제 (cst00309) artlSeq=5208397 file=extracted/1008/5208397/A형/2021년도 제32회 감정평가사 1차 2교시 A형.pdf
%% conversion: kordoc 4.24.51 -> tools/gen1.py (auto); review=proofread-sonnet; flags=none
\begin{question}{5}
국토의 계획 및 이용에 관한 법령상 용도지역·용도지구·용도구역에 관한 설명으로 옳지 않은 것은?
\begin{choices}
  \item 제2종일반주거지역은 중층주택을 중심으로 편리한 주거환경을 조성하기 위하여 필요한 지역을 말한다.
  \item 시·도지사는 대통령령으로 정하는 주거지역·공업지역·관리지역에 복합용도지구를 지정할 수 있다.
  \item 경관지구는 자연경관지구, 시가지경관지구, 특화경관지구로 세분할 수 있다.
  \item 관리지역에서 「농지법」에 따른 농업진흥지역으로 지정·고시된 지역은 농림지역으로 결정·고시된 것으로 본다.
  \item 시가화조정구역의 지정에 관한 도시·군관리계획의 결정은 시가화 유보기간이 끝난 날부터 그 효력을 잃는다.
\end{choices}
\end{question}
\answer{5}
